\begin{thebibliography}{22}

\bibitem{kister1957}
J. Kister, P. Stein, S. Ulam, W. Walden, and M. Wells, ``Experiments in chess,'' \emph{J. ACM}, vol. 4, no. 2, pp. 174--177, 1957.

\bibitem{mhalla2013}
M. Mhalla and F. Prost, ``Gardner's minichess variant is solved,'' \emph{ICGA J.}, vol. 36, no. 4, 2013, arXiv:1307.7118.

\bibitem{fairy}
Fairy-Stockfish developers, ``Fairy-Stockfish,'' GitHub repository. [Online]. Available: https://github.com/fairy-stockfish/Fairy-Stockfish

\bibitem{shannon1950}
C. E. Shannon, ``Programming a computer for playing chess,'' \emph{Philos. Mag.}, ser. 7, vol. 41, no. 314, pp. 256--275, 1950.

\bibitem{knuth1975}
D. E. Knuth and R. W. Moore, ``An analysis of alpha-beta pruning,'' \emph{Artif. Intell.}, vol. 6, no. 4, pp. 293--326, 1975.

\bibitem{samuel1959}
A. L. Samuel, ``Some studies in machine learning using the game of checkers,'' \emph{IBM J. Res. Develop.}, vol. 3, no. 3, pp. 210--229, 1959.

\bibitem{tesauro1995}
G. Tesauro, ``Temporal difference learning and TD-Gammon,'' \emph{Commun. ACM}, vol. 38, no. 3, pp. 58--68, 1995.

\bibitem{baxter2000}
J. Baxter, A. Tridgell, and L. Weaver, ``Learning to play chess using temporal differences,'' \emph{Mach. Learn.}, vol. 40, no. 3, pp. 243--263, 2000.

\bibitem{kocsis2006}
L. Kocsis and C. Szepesv\'ari, ``Bandit based Monte-Carlo planning,'' in \emph{Machine Learning: ECML 2006}, ser. Lecture Notes in Computer Science, vol. 4212. Berlin, Germany: Springer, 2006, pp. 282--293.

\bibitem{coulom2007}
R. Coulom, ``Efficient selectivity and backup operators in Monte-Carlo tree search,'' in \emph{Computers and Games (CG 2006)}, ser. Lecture Notes in Computer Science, vol. 4630. Berlin, Germany: Springer, 2007, pp. 72--83.

\bibitem{browne2012}
C. B. Browne \emph{et al.}, ``A survey of Monte Carlo tree search methods,'' \emph{IEEE Trans. Comput. Intell. AI Games}, vol. 4, no. 1, pp. 1--43, Mar. 2012.

\bibitem{silver2018}
D. Silver \emph{et al.}, ``A general reinforcement learning algorithm that masters chess, shogi, and Go through self-play,'' \emph{Science}, vol. 362, no. 6419, pp. 1140--1144, 2018.

\bibitem{silver2017}
D. Silver \emph{et al.}, ``Mastering the game of Go without human knowledge,'' \emph{Nature}, vol. 550, pp. 354--359, 2017.

\bibitem{sun2021}
M. Sun and R. Tan, ``Alpha-Mini: Minichess agent with deep reinforcement learning,'' 2021, arXiv:2112.13666.

\bibitem{hammersborg2022}
P. Hammersborg and I. Str\"umke, ``Reinforcement learning in an adaptable chess environment for detecting human-understandable concepts,'' 2022, arXiv:2211.05500.

\bibitem{koyamada2023}
S. Koyamada \emph{et al.}, ``Pgx: Hardware-accelerated parallel game simulators for reinforcement learning,'' in \emph{Adv. Neural Inf. Process. Syst. (NeurIPS), Datasets and Benchmarks Track}, 2023, arXiv:2303.17503.

\bibitem{sutton2018}
R. S. Sutton and A. G. Barto, \emph{Reinforcement Learning: An Introduction}, 2nd ed. Cambridge, MA, USA: MIT Press, 2018.

\bibitem{russell2021}
S. Russell and P. Norvig, \emph{Artificial Intelligence: A Modern Approach}, 4th ed. Hoboken, NJ, USA: Pearson, 2021.

\bibitem{zobrist1970}
A. L. Zobrist, ``A new hashing method with application for game playing,'' Comput. Sci. Dept., Univ. Wisconsin, Madison, WI, USA, Tech. Rep. 88, 1970.

\bibitem{rosin1997}
C. D. Rosin and R. K. Belew, ``New methods for competitive coevolution,'' \emph{Evol. Comput.}, vol. 5, no. 1, pp. 1--29, 1997.

\bibitem{kingma2015}
D. P. Kingma and J. Ba, ``Adam: A method for stochastic optimization,'' in \emph{Proc. Int. Conf. Learn. Represent. (ICLR)}, 2015, arXiv:1412.6980.

\bibitem{bradley1952}
R. A. Bradley and M. E. Terry, ``Rank analysis of incomplete block designs: I. The method of paired comparisons,'' \emph{Biometrika}, vol. 39, no. 3/4, pp. 324--345, 1952.

\end{thebibliography}